The holographic grating is a type of diffraction grating. The term holographic arises from the fact that the technique with which the diffraction grid is made consists in having a transparent film in which an interference figure is impressed. When a non-monochromatic light hits a diffraction grating at a distance $ z $, the outgoing light does not remain unchanged but is broken down into its chromatic components. Each point on the grating is struck by the light coming from the white source at a different angle. From each point then the light will be diffracted at different angles, as shown in the image (\ref{Diff}). The result is that these rays diverge as if they came from a point located in the plane of the source but shifted by $\Delta x$ and this happens for all colors. The distance between the virtual source and the real source is linked by the following equation:
%aggiungere immagine della difrazione dovuta dal reticolo olografico
\begin{equation}
    \Delta x= z\tan{\phi}\approx z\sin{\phi}=\frac{z\lambda}{d}
    \label{deltax}
\end{equation}
Where $z$ is the distance between the source and the grating, in fact farther the grating is from the source, greater is the angle that separates the virtual source and the real source; $\lambda$ is the wavelength of the virtual source and $d$ is the step of the holographic grating. 

During the experiment it is necessary to measure the wavelength through the holographic grating with the grating groove space $d$ known with the following equation:
\begin{equation}
    \lambda=\frac{\Delta x d}{z}
\end{equation}
And then, using the same configuration, replace the holographic lattice with the second one and get $ d '$:
\begin{equation}
    d'=\frac{z\lambda}{\Delta x'}
    \label{d}
\end{equation}
\begin{figure}[htp!]
  \centering
  \includegraphics[scale=0.25]{Trattazione teorica/dsa.png} 
  \caption{From this image it is possible to schematically observe the virtual sources produced by the diffraction.}
  \label{Diff}
\end{figure}